\section{Soft and collinear limits: evolution equations and further applications}
\label{sec:limits}

The trijet cross-section contains, as limiting cases, the first step of two different evolutions: the
high-energy (B--JIMWLK/BK) evolution of the target correlators, which is generated by soft gluons, and
the DGLAP evolution of the final-state fragmentation, which is generated by collinear gluons. For
massless quarks this has been shown in ref.~\cite{Iancu:2022gpw}. In this section we study how these
limits are modified by the quark mass. We shall need the leading-order cross-section for the
production of a massive quark--antiquark pair by a transverse photon, which in our conventions reads
\begin{align}
\frac{\dd\sigma^{\gT A\to\qqbar+X}_{\rm LO}}{\dd^3k_1\,\dd^3k_2}={}&\frac{2\,\aem\,\Nc}{(2\pi)^6\,q^+}\;
\delta\big(q^+-k_1^+-k_2^+\big)\sum_fe_f^2\int_{\x,\y,\bx,\by}e^{-i\kk_1\cdot(\x-\bx)-i\kk_2\cdot(\y-\by)}
\nonumber\\
&\times\Phi_{\rm LO}(\bm{r},\bar{\bm r})\;\cW_{\rm LO}(\x,\y;\bx,\by)\,,
\label{eq:LO}\\[1mm]
\Phi_{\rm LO}(\bm{r},\bar{\bm r})\equiv{}&\big[\zeta^2+(1-\zeta)^2\big]\,\frac{\bm{r}\cdot\bar{\bm r}}{r\,\bar r}\;
\ef^2\,K_1(\ef r)\,K_1(\ef\bar r)+\mf^2\,K_0(\ef r)\,K_0(\ef\bar r)\,,
\label{eq:PhiLO}
\end{align}
with $\bm r=\x-\y$, $\bar{\bm r}=\bx-\by$ and $\cW_{\rm LO}$ given in eq.~\eqref{eq:WLO}. It follows from
the vertex \eqref{eq:Hgamma2}, whose Fourier transform is the leading-order wave function
\begin{align}
\psi^i_{\lambda_1\lambda_2}(\bm r)&\equiv\varphi^{ij}_{\lambda_1\lambda_2}(\zeta)\,\frac{r^j}{r}\,\ef K_1(\ef r)
-i\,\mf\,\mu^i_{\lambda_1\lambda_2}\,K_0(\ef r)\,,
\nonumber\\
\sum_{\lambda_1,\lambda_2}\sum_i&\;\psi^i_{\lambda_1\lambda_2}(\bm r)\,\psi^{i\,*}_{\lambda_1\lambda_2}(\bar{\bm r})
=4\,\Phi_{\rm LO}(\bm r,\bar{\bm r})\,.
\label{eq:psiLO}
\end{align}

\subsection{The soft-gluon limit and the B--JIMWLK evolution}
\label{sec:soft}

\subsubsection{Factorisation of the soft gluon}

Consider the limit $\xi\to0$ at fixed transverse coordinates. The emitter does not recoil, $\yp\to\y$
and $\xp\to\x$, so that $\R,\Rt\to\bm r=\x-\y$; furthermore $\eta\to1-\zeta$, $\eft\to\ef$, and
\begin{equation}
a,\;\tilde a\to\frac{\xi}{\zeta(1-\zeta)}\,,\qquad\of\to\frac{\xi\,\mf}{1-\zeta}\,,\qquad
\oft\to\frac{\xi\,\mf}{\zeta}\,.
\end{equation}
The mass affects this limit in a very simple way. In the radial functions \eqref{eq:Gab} the integral
is dominated by $u\sim a\Y^2/\R^2\ll1$, where $A_u\to\ef^2$. Changing variable to $s=a\Y^2/(u\R^2)$ one
finds
\begin{align}
\cG^{(11)}&\to\frac{\ef^2}{2\zeta(1-\zeta)\,\Y^2}\int_0^\infty\frac{\dd s}{1+s}\,K_2\big(\ef r\sqrt{1+s}\big)
=\frac{\ef K_1(\ef r)}{\zeta(1-\zeta)\,r}\;\frac{1}{\Y^2}\,,
\nonumber\\
\cG^{(01)}&\to\frac{K_0(\ef r)}{\zeta(1-\zeta)}\;\frac{1}{\Y^2}\,,
\label{eq:G-soft}
\end{align}
which coincide with the limits of $\cA^{(11)}$ and $\cA^{(01)}$ for $\of Y\ll1$. The two paths with a
helicity flip at the gluon vertex are suppressed by $\xi\of\propto\xi^2$, and the instantaneous terms by
$\xi$. Moreover $\tau^{mn}_{\lambda\lambda_2}(\xi,\eta)\to2(1-\zeta)\,\delta^{mn}\delta_{\lambda\lambda_2}$
and $\tau^{nm}_{\lambda_1\lambda}(\xi,\zeta)\to2\zeta\,\delta^{nm}\delta_{\lambda_1\lambda}$: a soft gluon
does not change the helicity of its emitter. Hence both effective vertices \eqref{eq:M} reduce to the
product of the leading-order wave function \eqref{eq:psiLO} and of the eikonal emission kernel,
\begin{equation}
\cM^{im}_{1,2\;\lambda_1\lambda_2}\to\frac{2}{\zeta}\;\frac{Y^m}{\Y^2}\;\psi^i_{\lambda_1\lambda_2}(\bm r)\,,\qquad
\tM^{im}_{1,2\;\lambda_1\lambda_2}\to\frac{2}{1-\zeta}\;\frac{X^m}{\X^2}\;\psi^i_{\lambda_1\lambda_2}(\bm r)\,,
\end{equation}
and the $\qqg$ component of the outgoing state becomes
\begin{align}
\big|\gamma_T^i\big\rangle_{\qqg}\;\underset{\xi\to0}{\simeq}\;&-\frac{e\,e_f\,g\,q^+}{(2\pi)^4}\int_{\x,\y,\z}
\int\dd\zeta\int\frac{\dd\xi}{\sqrt\xi}\;\delta^{(2)}(\w-\cc)\;\psi^i_{\lambda_1\lambda_2}(\bm r)
\nonumber\\
&\times\bigg\{\frac{Y^m}{\Y^2}\Big[U^{ab}(\z)\,V(\x)\,t^b\,V^\dagger(\y)-V(\x)\,V^\dagger(\y)\,t^a\Big]
\nonumber\\
&\qquad-\frac{X^m}{\X^2}\Big[U^{ab}(\z)\,V(\x)\,t^b\,V^\dagger(\y)-t^a\,V(\x)\,V^\dagger(\y)\Big]\bigg\}_{\alpha\beta}
\nonumber\\
&\times\big|\bar q^\beta_{\lambda_2}(\y)\;g^a_m(\z)\;q^\alpha_{\lambda_1}(\x)\big\rangle\,.
\label{eq:state-soft}
\end{align}
This is the leading-order $\qqbar$ state dressed by the Weizs\"acker--Williams field of its two colour
charges: the gluon is emitted either before the collision, in which case it scatters with the adjoint
Wilson line $U(\z)$, or after it, in which case the colour generator acts on the outgoing fermion. The
structure of eq.~\eqref{eq:state-soft} is exactly the one of the soft-gluon light-cone wave function
which underlies the derivation of the JIMWLK Hamiltonian from the
LCWF~\cite{Kovner:2000pt,Lublinsky:2016meo,Radhakrishnan:2026}, with the valence charges being here
those of the quark and of the antiquark.

The same degeneracy is visible in the helicity sums. Taking $\xi\to0$ in eqs.~\eqref{eq:H-direct},
\eqref{eq:Ht-direct} and \eqref{eq:Hx}, and including the explicit prefactors $\zeta^2$, $\eta^2$ and
$\zeta\eta$ of the three channels, one finds
\begin{equation}
\frac{\zeta^2}{\xi}\,\cH\simeq\frac{16}{\xi}\,\frac{\Y\cdot\bY}{\Y^2\,\bY^2}\,\Phi_{\rm LO}\,,\qquad
\frac{\eta^2}{\xi}\,\tH\simeq\frac{16}{\xi}\,\frac{\X\cdot\bX}{\X^2\,\bX^2}\,\Phi_{\rm LO}\,,\qquad
\frac{\zeta\eta}{\xi}\,\cH^\times\simeq\frac{16}{\xi}\,\frac{\X\cdot\bY}{\X^2\,\bY^2}\,\Phi_{\rm LO}\,,
\label{eq:soft-weights}
\end{equation}
for any of the four combinations of radial functions. The three channels thus involve the same factor
$\Phi_{\rm LO}$, including its mass term, and differ only by the emission kernel. Since the four terms
of each channel are multiplied by the same kernel, their colour structures can be added. The result is
\begin{align}
\frac{\dd\sigma^{\gT A\to\qqg+X}}{\dd^3k_1\,\dd^3k_2\,\dd^3k_3}\bigg|_{\xi\to0}={}&
\frac{2\,\aem\,\Nc}{(2\pi)^6\,q^+}\;\delta\big(q^+-k_1^+-k_2^+\big)\sum_fe_f^2
\nonumber\\
&\times\int_{\x,\y,\bx,\by}e^{-i\kk_1\cdot(\x-\bx)-i\kk_2\cdot(\y-\by)}\;\Phi_{\rm LO}(\bm r,\bar{\bm r})
\nonumber\\
&\times\frac{4\,\as\,\CF}{(2\pi)^4\,\xi\,q^+}\int_{\z,\bz}e^{-i\kk_3\cdot(\z-\bz)}\;
\Delta_{\rm real}\cQ(\x,\y,\by,\bx;\z,\bz)\,,
\label{eq:xs-soft}
\end{align}
where we neglected $k_3^+$ in the $\delta$-function. Here, with the eikonal emission kernels
\begin{equation}
\mathcal{K}_{XX}\equiv\frac{\X\cdot\bX}{\X^2\,\bX^2}\,,\qquad
\mathcal{K}_{YY}\equiv\frac{\Y\cdot\bY}{\Y^2\,\bY^2}\,,\qquad
\mathcal{K}_{XY}\equiv\frac{\X\cdot\bY}{\X^2\,\bY^2}\,,\qquad
\mathcal{K}_{YX}\equiv\frac{\Y\cdot\bX}{\Y^2\,\bX^2}\,,
\label{eq:kernels}
\end{equation}
and at large $\Nc$,
\begin{align}
\Delta_{\rm real}\cQ(\x,\y,\by,\bx;\z,\bz)&
\nonumber\\
\equiv{}&\left[\mathcal{K}_{XX}+\mathcal{K}_{YY}-\mathcal{K}_{XY}-\mathcal{K}_{YX}\right]\cQ(\x,\z,\bz,\bx)\,\cQ(\z,\y,\by,\bz)
\nonumber\\
&+\left[\mathcal{K}_{XX}+\mathcal{K}_{YY}\right]\cQ(\x,\y,\by,\bx)
-\left[\mathcal{K}_{XY}+\mathcal{K}_{YX}\right]S(\x,\y)\,S(\by,\bx)
\nonumber\\
&-\left[\mathcal{K}_{XX}-\mathcal{K}_{XY}\right]\cQ(\x,\y,\by,\bz)\,S(\bz,\bx)
-\left[\mathcal{K}_{YY}-\mathcal{K}_{YX}\right]\cQ(\x,\y,\bz,\bx)\,S(\by,\bz)
\nonumber\\
&-\left[\mathcal{K}_{XX}-\mathcal{K}_{YX}\right]S(\x,\z)\,\cQ(\z,\y,\by,\bx)
-\left[\mathcal{K}_{YY}-\mathcal{K}_{XY}\right]\cQ(\x,\z,\by,\bx)\,S(\z,\y)\,,
\label{eq:DeltaQ}
\end{align}
with $\X=\x-\z$, $\Y=\y-\z$ and similarly for the barred quantities. Equation~\eqref{eq:DeltaQ} is the
same as in the massless case~\cite{Iancu:2022gpw}: it is the action on the leading-order quadrupole of
the ``production'' version of the JIMWLK
Hamiltonian~\cite{Kovner:2006ge,Kovner:2006wr,Iancu:2013uva}, which describes the emission of a soft
gluon that is measured in the final state. The prefactor in the last line of
eq.~\eqref{eq:xs-soft} is fixed by the normalisation of the trijet cross-section and provides a check
of the latter: for a single emitter and without scattering it gives the familiar soft-gluon spectrum
$\dd N/(\dd k_3^+\dd^2\kk_3)=\as\CF/(\pi^2k_3^+\kk_3^2)$.

\subsubsection{Evolution equations}

If the soft gluon is not measured, the integration over $\kk_3$ sets $\bz=\z$ and the integration over
$k_3^+$ produces the logarithmic measure $\dd\xi/\xi$. Using $2\as\CF\simeq\as\Nc\equiv\pi\bar\alpha_s$,
eq.~\eqref{eq:xs-soft} then amounts to the following change of the colour structure of the
leading-order cross-section \eqref{eq:LO} under one step of the high-energy evolution,
\begin{equation}
\frac{\partial\,\cQ(\x,\y,\by,\bx)}{\partial\ln(1/\xi_0)}\bigg|_{\rm real}
=\frac{\bar\alpha_s}{2\pi}\int_{\z}\Delta_{\rm real}\cQ(\x,\y,\by,\bx;\z,\bz=\z)\,,
\label{eq:JIMWLK}
\end{equation}
where $\xi_0$ is the lower limit of the integration over the gluon fraction, so that
$\ln(1/\xi_0)$ is the rapidity interval available to the gluon. This is the set of terms of the B--JIMWLK equation for the
quadrupole~\cite{Dominguez:2011gc,Iancu:2011ns,Iancu:2011nj} in which the gluon is exchanged between a
fermion in the DA and a fermion in the CCA. We thus reach the following conclusions about which
evolution equations follow from the trijet cross-section.
\begin{itemize}
\item \emph{Production JIMWLK and the real part of the quadrupole evolution.} They follow directly from
eqs.~\eqref{eq:xs-soft}--\eqref{eq:JIMWLK}, with no further input. The emission kernel is the massless
one: the mass enters only through $\Phi_{\rm LO}$, i.e.\ through the impact factor. The leading
logarithmic high-energy evolution is therefore insensitive to the mass of the quarks of the projectile,
as expected since it is a property of the target. Quark masses affect the evolution kernel only at
next-to-leading order, through quark loops~\cite{Dai:2022imf}.
\item \emph{The remaining terms of the B--JIMWLK equation.} The terms in which the gluon is emitted and
absorbed on the same side of the cut, and the complete evolution of the dipoles $S(\x,\y)$ and
$S(\by,\bx)$ which appear in $\cW_{\rm LO}$, are generated by the virtual corrections, which are not part
of the trijet cross-section. They can however be reconstructed from the soft state
\eqref{eq:state-soft} by unitarity, as explained below.
\item \emph{The BK equation.} The BK equation is the evolution equation for the dipole $S$-matrix
which controls the total cross-section. It is obtained from the same light-cone wave function through
the optical theorem: the total cross-section is the sum over Fock states of the probability density of
each state times $2\,[1-\re S]$, with $S$ the forward $S$-matrix of that state. Only the component of
the $\qqg$ state before the collision is needed, i.e.\ the vertices $\cM_1$ and $\tM_1$. In the soft
limit eq.~\eqref{eq:state-soft} gives for the probability of the $\qqg$ state
\begin{equation}
\big|\psi_{\qqg}\big|^2\simeq\big|\psi_{\qqbar}\big|^2\;\frac{\as\CF}{\pi^2}\,\frac{\dd\xi}{\xi}\,\dd^2\z
\left|\frac{\X}{\X^2}-\frac{\Y}{\Y^2}\right|^2
=\big|\psi_{\qqbar}\big|^2\;\frac{\bar\alpha_s}{2\pi}\,\frac{\dd\xi}{\xi}\,\dd^2\z\;
\frac{(\x-\y)^2}{(\x-\z)^2(\z-\y)^2}\,,
\end{equation}
and the conservation of probability fixes the corresponding decrease of the probability of the
$\qqbar$ state. With $\cS_{\qqg}\simeq S(\x,\z)S(\z,\y)$ at large $\Nc$, one obtains the BK equation,
\begin{equation}
\frac{\partial S(\x,\y)}{\partial\ln(1/\xi_0)}=\frac{\bar\alpha_s}{2\pi}\int_{\z}
\frac{(\x-\y)^2}{(\x-\z)^2(\z-\y)^2}\,\Big[S(\x,\z)\,S(\z,\y)-S(\x,\y)\Big]\,,
\label{eq:BK}
\end{equation}
again with a mass-independent kernel, the mass dependence being entirely carried by the impact factor
$\int\dd\zeta\,\Phi_{\rm LO}(\bm r,\bm r)$ of the heavy-quark structure function.
\end{itemize}

Two comments on the range of validity of the soft approximation are in order. First,
eq.~\eqref{eq:G-soft} requires $a\Y^2\ll\R^2$, i.e.\ $\xi\,\Y^2\ll\zeta(1-\zeta)\,\bm r^2$: the lifetime of
the gluon fluctuation must be shorter than that of the parent pair. For larger gluon separations the
exact functions \eqref{eq:Gab} are cut off by the virtuality, whereas the eikonal kernel is not. This is
the origin of the kinematical (time-ordering) constraint which must be imposed on the high-energy
evolution beyond leading logarithmic
accuracy~\cite{Beuf:2014uia,Iancu:2015vea,Ducloue:2019ezk}; our exact expressions contain this
constraint, now with its dependence on the quark mass through $A_u$. Second, for a gluon emitted after
the collision the kernel $Y^m/\Y^2$ holds for $Y\ll1/\of=(1-\zeta)/(\xi\mf)$ only; at larger
distances the dead cone sets in. This affects the region of very large daughter dipoles, where the
integrand of the evolution equations is anyway suppressed by colour transparency, and it is a
correction of higher order in $\xi$ to the evolution.

\subsection{The collinear limit: massive splitting function and dead cone}
\label{sec:coll}

Consider now the limit in which the gluon is nearly collinear with the antiquark. As in the massless
case, we change the momentum variables of the antiquark--gluon pair to their total and relative
momenta,
\begin{equation}
\KK\equiv\kk_2+\kk_3\,,\qquad\PP\equiv\frac{\eta\kk_3-\xi\kk_2}{1-\zeta}\qquad\Longrightarrow\qquad
e^{-i\kk_2\cdot\y-i\kk_3\cdot\z}=e^{-i\KK\cdot\yp+i\PP\cdot\Y}\,,
\end{equation}
so that $\PP$ is conjugate to the separation $\Y$ between the gluon and the antiquark and coincides
with the momentum $\pt$ of eq.~\eqref{eq:relmom}; the splitting fraction of the gluon is
\begin{equation}
z\equiv\frac{\xi}{1-\zeta}=\frac{k_3^+}{k_2^++k_3^+}\,,\qquad\of=z\,\mf\,.
\end{equation}
The collinear region is $|\PP|\ll|\kk_1|,|\KK|$, i.e.\ large $Y$. A gluon emitted before the collision
cannot be at a large distance from its emitter, because of the virtuality barrier in the functions
$\cG^{(\alpha\beta)}$ and $\cG_E$, and the interference terms are suppressed because the gluon cannot be
collinear to both fermions. The collinear limit is therefore controlled by the last term of
eq.~\eqref{eq:xs-qbar}, in which the gluon is emitted after the collision in the DA and in the CCA.
Its colour structure, eq.~\eqref{eq:W-qbar}, is the leading-order one, $\cW_{\rm LO}(\x,\yp;\bx,\byp)$,
and depends on $\x$, $\yp$, $\bx$ and $\byp$ only. Changing the integration variables from $(\y,\z)$ to
$(\yp,\Y)$, the integrals over $\Y$ and $\bY$ factorise and are computed with eq.~\eqref{eq:FT-fact}:
\begin{equation}
\int_{\Y}e^{i\PP\cdot\Y}\,\frac{Y^n}{Y}\,\of K_1(\of Y)=\frac{2\pi i\,P^n}{\PP^2+\of^2}\,,\qquad
\int_{\Y}e^{i\PP\cdot\Y}\,K_0(\of Y)=\frac{2\pi}{\PP^2+\of^2}\,.
\end{equation}
Inserting the factorised functions \eqref{eq:Aexplicit} in eq.~\eqref{eq:Hrr}, the term proportional to
$(\Y\times\bY)$ vanishes after the integration, since it becomes $\PP\times\PP=0$, and the four terms in
which two different helicity paths interfere cancel pairwise. We find
\begin{equation}
\int_{\Y,\bY}e^{i\PP\cdot(\Y-\bY)}\;\cH[\cA;\bar\cA]
=\frac{8\,(2\pi)^2}{\zeta^2(1-\zeta)^2}\;\Phi_{\rm LO}(\R,\bR)\;
\frac{\big[\eta^2+(1-\zeta)^2\big]\PP^2+\xi^2\of^2}{\big(\PP^2+\of^2\big)^2}\,,
\label{eq:coll-H}
\end{equation}
that is, the helicity sum factorises exactly into the leading-order photon factor \eqref{eq:PhiLO}, with
its mass term, and a factor which describes the emission of the gluon by the massive antiquark. The
latter can be written as
\begin{equation}
\frac{\big[\eta^2+(1-\zeta)^2\big]\PP^2+\xi^2\of^2}{\big(\PP^2+\of^2\big)^2}
=\frac{2z\,(1-\zeta)^2}{\PP^2+z^2\mf^2}\;\cP_{q\to g}\big(z,\PP^2\big)\,,
\end{equation}
with
\begin{equation}
\boxed{\;\cP_{q\to g}\big(z,\PP^2\big)\equiv\frac{1+(1-z)^2}{2z}-\frac{z(1-z)\,\mf^2}{\PP^2+z^2\mf^2}\;}
\label{eq:Pmassive}
\end{equation}
the quasi-collinear splitting function of a massive quark~\cite{Catani:2000ef}, in the normalisation in
which the massless one is $P_{q\to g}(z)=[1+(1-z)^2]/2z$. Indeed, the squared invariant mass of the
antiquark--gluon pair is $(k_2+k_3)^2-\mf^2=(\PP^2+z^2\mf^2)/[z(1-z)]$, so that the second term of
eq.~\eqref{eq:Pmassive} is the familiar $-\mf^2/[(k_2+k_3)^2-\mf^2]$. Collecting the prefactors, the
trijet cross-section in the collinear region reads
\begin{equation}
\frac{\dd\sigma^{\gT A\to\qqg+X}}{\dd^3k_1\,\dd^3k_2\,\dd^3k_3}\bigg|_{\rm coll}\simeq
\frac{\dd\sigma^{\gT A\to\qqbar+X}_{\rm LO}}{\dd^3k_1\,\dd^3p_2}\bigg|_{p_2^+=k_2^++k_3^+,\;\pp_2=\KK}\times
\frac{\as\CF}{\pi^2}\;\frac{1}{(1-\zeta)\,q^+}\;\frac{\cP_{q\to g}\big(z,\PP^2\big)}{\PP^2+z^2\mf^2}\,.
\label{eq:xs-coll}
\end{equation}
At fixed momentum of the parent antiquark, $\dd k_3^+\,\dd^2\kk_3=(1-\zeta)q^+\,\dd z\,\dd^2\PP$, hence
eq.~\eqref{eq:xs-coll} is the leading-order cross-section times the emission probability
\begin{equation}
\dd\mathcal{N}=\frac{\as\CF}{2\pi}\;\dd z\;\frac{\dd\PP^2}{\PP^2+z^2\mf^2}
\left[\frac{1+(1-z)^2}{z}-\frac{2z(1-z)\,\mf^2}{\PP^2+z^2\mf^2}\right]\,.
\label{eq:dN}
\end{equation}
The corresponding result for an emission by the quark follows from $\zeta\leftrightarrow\eta$.

Equations~\eqref{eq:xs-coll} and \eqref{eq:dN} exhibit the dead-cone
effect~\cite{Dokshitzer:1991fd}: the emission probability behaves as $\PP^2/(\PP^2+z^2\mf^2)^2$ for
$|\PP|\ll z\mf$, so that the radiation is suppressed at angles
$\theta\simeq|\PP|/[z(1-z)p_2^+]$ smaller than $\mf/[(1-z)p_2^+]$, i.e.\ smaller than the ratio of the
mass to the energy of the heavy quark (up to the normalisation of the light-cone momenta). In transverse coordinates this is the exponential decay of the
kernel $\of K_1(\of Y)/Y$ for $Y\gtrsim1/\of$ which we noticed in section~\ref{sec:after}. As a
consequence, the integral over the relative momentum is finite,
\begin{equation}
\int_0^{\mu^2}\dd\PP^2\;\frac{2\,\cP_{q\to g}(z,\PP^2)}{\PP^2+z^2\mf^2}
=\frac{1+(1-z)^2}{z}\,\ln\frac{\mu^2}{z^2\mf^2}-\frac{2(1-z)}{z}+\order{\frac{z^2\mf^2}{\mu^2}}\,.
\label{eq:coll-int}
\end{equation}
There is no collinear divergence to be absorbed in a bare fragmentation function, in contrast with the
massless case, where the same integral requires an infrared cutoff $\Lambda$ and yields
$\ln(\mu^2/\Lambda^2)$~\cite{Iancu:2022gpw}. The logarithm $\ln(\mu^2/\mf^2)$ of
eq.~\eqref{eq:coll-int}, which becomes large when the transverse momenta of the jets are much larger
than the quark mass, is the first term of the series which is resummed by the DGLAP
evolution~\cite{Gribov:1972ri,Altarelli:1977zs,Dokshitzer:1977sg} of the heavy-quark fragmentation
function between the scales $\mf$ and $\mu$, and the non-logarithmic terms contribute to its
perturbative initial condition~\cite{Mele:1990cw}; the constant depends on the way in which the
integral is cut off (here, a sharp cutoff on $\PP^2$). The massless limit of eq.~\eqref{eq:coll-H} is
$16\,(2\pi)^2\,\Phi_{\rm LO}\;z\,P_{q\to g}(z)/[\zeta^2\PP^2]$, in agreement with eq.~(7.3) of
ref.~\cite{Iancu:2022gpw}; it is approached only for $|\PP|\gg z\mf$, which makes precise the
non-uniformity of the massless limit discussed in section~\ref{sec:massless}.

\subsection{Further applications}
\label{sec:apps}

We close this section with a list of other quantities which can be obtained from the results of this
paper.
\begin{enumerate}
\item \emph{Real NLO corrections to heavy-quark dijet and dihadron production.} Integrating the trijet
cross-section over the momentum of the gluon gives the real part of the NLO corrections to
$\gT A\to Q\bar Q+X$. The integral over $\kk_3$ sets $\bz=\z$ and gives a factor $(2\pi)^2$, and the
integral over $k_3^+$ removes the $\delta$-function. The colour structure simplifies when the three
partons scatter on both sides of the cut,
$\cW(\x,\y,\z;\bx,\by,\z)\simeq S(\x,\bx)\,S(\by,\y)-S(\x,\z)S(\z,\y)-S(\z,\bx)S(\by,\z)+1$. Integrating
instead over the quark or the antiquark gives the quark--gluon and antiquark--gluon dijet
cross-sections, and integrating over two partons the single-inclusive ones. Together with the
virtual corrections, the soft and collinear logarithms identified above are then to be absorbed into
the B--JIMWLK evolution of the target and into the evolution of the fragmentation functions or the
jet definition.
\item \emph{Heavy-quark structure functions.} Through the optical theorem, the vertices $\cM_1$ and
$\tM_1$ give the $\qqg$ contribution to the transverse structure function with massive quarks, to be
compared with refs.~\cite{Beuf:2021srj,Beuf:2022ndu}.
\item \emph{Diffraction.} The coherent cross-section of section~\ref{sec:mom} is the diffractive trijet
cross-section. Integrated over the final state at fixed diffractive mass, it gives the $\qqg$
contribution to the diffractive structure functions with heavy quarks, which dominates at large
diffractive masses~\cite{Wusthoff:1997fz,GolecBiernat:1999qd,Kovchegov:1999ji,Beuf:2022kyp}, and in the
regime where the gluon is soft and at large distance it gives access to the gluon dipole and to
diffractive jets~\cite{Iancu:2021rup,Iancu:2022lcw}.
\item \emph{Photoproduction.} All our formulas have a finite limit for $Q^2\to0$, where $\ef\to\mf$,
$A_u\to\mf^2(1+u\zeta\xi/\eta)$ and $\Ef^2\to\mf^2(1/\zeta+1/\eta)$. They can therefore be used for the
photoproduction of heavy-quark trijets in ultra-peripheral collisions, where the photon is real.
\item \emph{Back-to-back and TMD limits.} When the transverse momentum imbalance of the three jets is
much smaller than their individual momenta, the Wilson-line correlators can be expanded around
coincident points and the cross-section factorises into hard factors and gluon
TMDs~\cite{Dominguez:2011wm,Altinoluk:2020qet,Altinoluk:2021ygv}. At leading order the quark mass is known to make the linearly polarised gluon
distribution contribute even for real photons~\cite{Marquet:2017xwy}; the corresponding hard factors
for three jets with massive quarks can be extracted from our results. For back-to-back heavy-quark pairs, the soft gluons of
eq.~\eqref{eq:state-soft} emitted outside the jets generate the Sudakov double
logarithms~\cite{Mueller:2013wwa,Caucal:2022ulg}.
\item \emph{Dilute limit.} Expanding the Wilson lines to second order in the target field gives the
impact factor for $\gT\to Q\bar Qg$ in $k_\perp$-factorisation, which is the starting point for a
comparison with fixed-order calculations in collinear factorisation and for the matching to the BFKL
regime.
\end{enumerate}
